\subsubsection{Vaizdų duomenų aibė – Muffin vs Chihuahua}
\label{subsubsec:muffin_vs_chihuahua}

Vaizdų duomenų aibė paimta iš Kaggle platformos \cite{muffin_vs_chihuahua}. Aibę sudaro apie $5\,917$ JPG formato spalvotų (RGB) nuotraukų, padalintų į dvi klases: čihuahua šuo (angl. \textit{chihuahua}) ir keksiukas (angl. \textit{muffin}). Duomenys pirminiame šaltinyje yra suskaidyti į dvi pagrindines aplankų grupes – mokymo (angl. \textit{train}) ir testavimo (angl. \textit{test}), o kiekvienoje iš jų yra po du poaplankius pagal klases (\textit{chihuahua} ir \textit{muffin}).

Iš viso aibėje yra apie $3\,199$ čihuahua ir $2\,718$ keksiukų atvaizdų, t.~y. klasių disbalanso santykis $\rho \approx 1{,}22$, kuris laikomas pakankamai subalansuotu. Vaizdai yra skirtingų matmenų ir proporcijų – juose pasitaiko ir spalvotų (RGB), ir nespalvotų (grayscale) nuotraukų, kai kurios turi baltus rėmelius ar foną. Apibendrinta informacija apie aibę pateikta \ref{tab:muffin_vs_chihuahua_aprasymas} lentelėje.

\begin{table}[H]
    \centering
    \caption{Vaizdų duomenų aibės „Muffin vs Chihuahua“ aprašymas.}
    \begin{tabular}{|l|l|}
        \hline
        Charakteristika & Reikšmė \\
        \hline
        Šaltinis & Kaggle (Samuel Cortinhas) \\ \hline
        Užduoties tipas & Dvinarė vaizdų klasifikacija \\ \hline
        Klasių skaičius & $2$ (\textit{chihuahua}, \textit{muffin}) \\ \hline
        Bendras vaizdų skaičius & $\approx 5\,917$ \\ \hline
        Klasės \textit{chihuahua} vaizdų & $\approx 3\,199$ \\ \hline
        Klasės \textit{muffin} vaizdų & $\approx 2\,718$ \\ \hline
        Disbalanso santykis $\rho$ & $\approx 1{,}22$ \\ \hline
        Failo formatas & JPG (RGB, dalis – grayscale) \\ \hline
        Vaizdų matmenys & Skirtingi (nestandartizuoti) \\ \hline
        Pradinis padalijimas & \texttt{train} ir \texttt{test} aplankai \\ \hline
    \end{tabular}
    \label{tab:muffin_vs_chihuahua_aprasymas}
\end{table}